\documentclass[12pt,a4paper]{article}
\usepackage[margin=25mm, headheight=15pt, headsep=10pt]{geometry}
\usepackage{fontspec}
\usepackage[table]{xcolor} % Note: table option is required for rowcolors
\usepackage{titlesec}
\usepackage{tocloft}
\usepackage{fancyhdr}
\usepackage{booktabs}
\usepackage{tcolorbox}
\tcbuselibrary{skins, breakable}
\usepackage[colorlinks=true, linkcolor=emerald, urlcolor=emerald, bookmarks=true]{hyperref}

\definecolor{emerald}{RGB}{0,102,51}
\definecolor{darkred}{RGB}{139,0,0}
\definecolor{lightgray}{RGB}{245,245,245}

\usepackage[nil]{babel}
\babelprovide[import, main]{bengali}
\babelprovide[import]{arabic}
\babelfont{rm}[Renderer=HarfBuzz,Scale=1.0]{Noto Serif Bengali}
\babelfont{sf}[Renderer=HarfBuzz]{Noto Sans Bengali}
\babelfont{tt}[Renderer=HarfBuzz]{Noto Sans Bengali}
\babelfont[arabic]{rm}[Renderer=HarfBuzz,Scale=1.1]{Noto Naskh Arabic}

% Section styling
\titleformat{\section}{\Large\bfseries\color{emerald}}{\thesection.}{0.5em}{}
\titleformat{\subsection}{\large\bfseries\color{emerald!85!black}}{\thesubsection.}{0.5em}{}
\titleformat{\subsubsection}{\normalsize\bfseries\color{emerald!70!black}}{\thesubsubsection.}{0.5em}{}
\titlespacing*{\section}{0pt}{18pt}{8pt}

% Fancy Header/Footer setup
\pagestyle{fancy}
\fancyhf{} % clear all header and footer fields
\fancyhead[L]{\color{emerald}\small\textbf{\nouppercase{\leftmark}}}
\fancyfoot[L]{\scriptsize\color{black!50}{\lat{AI}}-সংকলিত, আলেম-যাচাইকৃত নয়}
\fancyfoot[C]{\color{emerald!70!black}\thepage}
\renewcommand{\headrulewidth}{0.4pt}
\renewcommand{\headrule}{\hbox to\headwidth{\color{emerald}\leaders\hrule height \headrulewidth\hfill}}
\renewcommand{\footrulewidth}{0pt}

% Custom boxes
% 1. Ayat/Hadith Box
\newtcolorbox{ayatbox}[1][]{
    enhanced, breakable, 
    colback=emerald!5, 
    colframe=emerald, 
    boxrule=0.5pt, 
    arc=3pt, 
    left=10pt, right=10pt, top=10pt, bottom=10pt,
    #1
}

% 2. Quote/Translation Box
\newtcolorbox{quotebox}[1][]{
    enhanced, breakable,
    frame hidden,
    colback=lightgray,
    borderline west={3pt}{0pt}{emerald},
    left=15pt, right=10pt, top=10pt, bottom=10pt,
    sharp corners,
    #1
}

% 3. AI-generated notice box
\newtcolorbox{warnbox}[1][]{
    enhanced, breakable,
    colback=red!4,
    colframe=darkred,
    boxrule=0.8pt,
    arc=3pt,
    left=10pt, right=10pt, top=10pt, bottom=10pt,
    #1
}

% Commands
\newcommand{\arab}[1]{\foreignlanguage{arabic}{#1}}
\newcommand{\kib}[1]{\textcolor{emerald}{\textbf{#1}}}

% Latin (English) fallback: Noto Serif Bengali has NO Latin glyphs
\newfontfamily\latfont[Renderer=HarfBuzz]{Noto Serif}
\newcommand{\lat}[1]{{\latfont #1}}

\setlength{\parskip}{5pt}
\setlength{\parindent}{0pt}

% Fix for missing bullet in Noto Serif Bengali
\renewcommand{\labelitemi}{$\bullet$}



\begin{document}
\begin{center}{\Large\bfseries\color{emerald} পার্ট ৫খ — মাসআলা: সাহু সিজদা, ভুল ও সংশোধন, সন্দেহ}\end{center}
\vspace{1em}
\section*{পার্ট ৫খ — মাসআলা: সাহু সিজদা, ভুল ও সংশোধন, সন্দেহ}

\begin{warnbox}
 \textbf{গুরুত্বপূর্ণ নোটিশ (\lat{Important} \lat{Notice}):} এই কন্টেন্টটি \lat{AI} (কৃত্রিম বুদ্ধিমত্তা) দিয়ে প্রস্তুত ও সংকলিত। \lat{AI} কঠোর সূত্র-মান নীতি মেনে শুধু নির্ভরযোগ্য উৎস থেকে তথ্য সংগ্রহ করেছে — কেবল সহীহ/হাসান হাদিস ও সহীহ সনদের আসার রাখা হয়েছে, দুর্বল সনদ বাদ দেওয়া হয়েছে। তবে এটি কোনো আলেম বা মুহাদ্দিস কর্তৃক যাচাইকৃত নয়।
\end{warnbox}

\begin{quote}
\textbf{সম্পর্কিত ফাইল:} পার্ট ৩ (শ্রেণিবিভাগ), পার্ট ৪ (পদ্ধতি), পার্ট ৫ক (জামাত-মাসবুক), পার্ট ৬ (মাযহাব-তুলনা), পার্ট ৭খ (ভুল-ও-প্রতিকার)।
\end{quote}

\begin{quote}
\textbf{যে প্রশ্নের উত্তর এই অংশে:} নামাজে ভুল করলে কী হয়? রাকাত কম/বেশি পড়লে কী করবেন? সন্দেহ হলে কীভাবে মীমাংসা? সাহু সিজদা কখন, কীভাবে, কখন নয়?
\end{quote}

\medskip\hrule\medskip

\subsection*{১. সাহু সিজদা কী ও কেন}

\textbf{সাহু সিজদা} = নামাজের \textbf{ভুল} হওয়ায় শেষে (সালামের আগে বা পরে) দুটি সিজদা — নামাজের ত্রুটি পূরণের জন্য।

\textbf{সহীহ হাদিস — ভিত্তি (সহীহ বুখারি ৪০১; সহীহ মুসলিম ৫৭৩):}

\begin{quote}
আবু হুরাইরা (রাঃ) থেকে: নবীজি (সা.) পাঁচ রাকাত নামাজ পড়লেন (অর্থাৎ ভুল করে)। সালামের পর লোকেরা জিজ্ঞেস করল: নামাজ কি বাড়ানো হয়েছে? নবীজি (সা.) বললেন: \textbf{\lat{“}তোমাদের প্রত্যেকের সন্দেহকে আমি (ভুলকে) সাহু সিজদা দিয়ে ঠিক করি।\lat{”}} — এরপর নবীজি (সা.) \textbf{সালামের পর দুই সিজদা} করলেন।
\end{quote}

\textbf{সহীহ হাদিস — মূল নীতি (সহীহ মুসলিম ৫৭২):}

\begin{quote}
আবদুল্লাহ ইবনে মাসউদ (রাঃ) থেকে: নবীজি (সা.) \textbf{সালামের আগে} সাহু সিজদা করেছেন (মাঝে ৩ রাকাত-সংশোধনের ঘটনা)।
\end{quote}

\begin{quote}
\textbf{সারমর্ম:} নবীজি (সা.) কখনো \textbf{সালামের আগে}, কখনো \textbf{পরে} সাহু সিজদা করেছেন — পরিস্থিতি অনুযায়ী। ফকিহরা এই দুই পথের নিয়ম মাযহাবভেদে সাজিয়েছেন (নিচে ৪ নম্বর অংশ)।
\end{quote}

\medskip\hrule\medskip

\subsection*{২. ভুল কী কী হতে পারে — দ্রুত সারণি}

\begin{table}[h]\centering\small
\begin{tabular}{lll}
ভুলের ধরন & উদাহরণ & কী করবেন (সংক্ষেপ) \\
\hline
\textbf{রুকন ছুটে গেছে} (যেমন সিজদা, রুকু, কিরাআত-এর অংশ) & ভুলে সিজদা না করা & যদি মনে পড়ে এক রুকনে যাওয়ার আগে — ফিরে গিয়ে করবে; পরে মনে পড়লে পুরো রাকাত পুনরায় + সাহু সিজদা (বিস্তারিত পার্ট ৭খ) \\
\textbf{ওয়াজিব ছুটে গেছে} (হানাফি পরিভাষা) & তাশাহহুদ, সূরা-মিলান, সালাম & সাহু সিজদা \\
\textbf{রাকাত কম/বেশি} & ৪-এর বদলে ৩ বা ৫ & সন্দেহ দূর করে নিশ্চিত সংখ্যায় মীমাংসা + সাহু সিজদা \\
\textbf{সন্দেহ} (কত রাকাত হলো?) & ৩ না ৪ — অনিশ্চিত & নিশ্চিত ন্যূনতম ধরে (৩) নেওয়া + সাহু সিজদা (হানাফি); জমহুরের কয়েকটি দৃষ্টিভঙ্গি (নিচে ৫) \\
\textbf{সালামের আগে/পরে ভুল-সংশোধন} & সাহু সিজদার সময় & মাযহাবভেদে — নিচে ৪ \\
\hline
\end{tabular}
\end{table}

\medskip\hrule\medskip

\subsection*{৩. সাহু সিজদার পদ্ধতি (সহীহ হাদিসের রূপরেখা)}

\begin{enumerate}
\item \textbf{সালামের আগে:} শেষ তাশাহহুদ পড়ে $\rightarrow$ ডান-বামে সালাম না ফিরে $\rightarrow$ \textbf{দুটি সিজদা} $\rightarrow$ আবার তাশাহহুদ (হানাফিতে দুরুদ-দুআসহ) $\rightarrow$ সালাম।
\item \textbf{সালামের পরে:} সালাম ফিরানোর পর \textbf{দুটি সিজদা} $\rightarrow$ আবার তাশাহহুদ (হানাফি) $\rightarrow$ সালাম।
\end{enumerate}
\begin{quote}
\textbf{কোনটি কখন?} সাধারণ নিয়ম (হানাফি-প্রধান): \\
- ভুল যদি \textbf{নামাজের ভেতরে} (তাশাহহুদ ছাড়া ইত্যাদি) — \textbf{সালামের আগে}। \\
- ভুল যদি \textbf{সালামের পরে} (সালাম-সংক্রান্ত, যেমন ভুলে ২ সালাম) — \textbf{সালামের পরে}। \\
- শাফেয়ি-হাম্বলি মতেও দুই অবস্থানই আছে — প্রধান নিয়ম "সালামের আগে" (আরও বিস্তারিত পার্ট ৬)।
\end{quote}

\medskip\hrule\medskip

\subsection*{৪. রাকাত কম/বেশি পড়লে — মাযহাবগত মীমাংসা}

\subsubsection*{৪.১ হানাফি মাযহাব}

\begin{itemize}
\item \textbf{কম পড়লেন} (নিশ্চিত): পুরো রাকাত যোগ করে শেষে সাহু সিজদা।
\item \textbf{বেশি পড়লেন} (নিশ্চিত): নামাজ শেষের দিকে সাহু সিজদা দিয়ে নামাজ \textbf{বাতিল নয়} — সাহু সিজদা দিলে ঠিক হয় (বুখারি ৪০১-এর ঘটনা)।
\item \textbf{সন্দেহ:} যতটুকু \textbf{নিশ্চিত} (ন্যূনতম) সে রাকাতে ধরবেন — যেমন ৩/৪ সন্দেহ $\rightarrow$ ৩ ধরে নেওয়া $\rightarrow$ শেষে সাহু সিজদা।
\end{itemize}
\subsubsection*{৪.২ শাফেয়ি-মালেকি-হাম্বলি (জমহুর)}

\begin{itemize}
\item \textbf{সন্দেহে সাধারণ নিয়ম:} হাদিসে এসেছে — \textbf{\lat{“}তোমাদের কেউ নামাজে সন্দেহ করলে সে যেন সন্দেহের প্রস্রবণ (বাতিল) করে নিশ্চিততার দিকে ঝোঁকে, তারপর সাহু সিজদা করে।\lat{”}} (সহীহ মুসলিম ৫৭২)।
\item প্রসিদ্ধ মত: সন্দেহ হলে \textbf{নিশ্চিত ন্যূনতম} ধরুন + সাহু সিজদা (হানাফির মতো) — অথবা কিছু বর্ণনায় \textbf{সবচেয়ে সঠিক-অনুমান} নেওয়া যায় (বিস্তারিত পার্ট ৬)।
\end{itemize}
\begin{quote}
\textbf{ব্যবহারিক দিক:} চার মাযহাবের মূল এক জিনিস — \textbf{সন্দেহ হলে নিশ্চিত দিকে যান, তারপর সাহু সিজদা দিয়ে মীমাংসা করুন।} কাতারের সন্দেহ নিয়ে নামাজ ছেড়ে দেওয়া বা পুনরায় শুরু করা \textbf{সুন্নাহ-বিরুদ্ধ}।
\end{quote}

\medskip\hrule\medskip

\subsection*{৫. সাহু সিজদা কখন \textbf{না} দিলেই চলে}

\begin{itemize}
\item \textbf{সুন্নত/মুস্তাহাব} ভুলে ছুটে গেলে সাহু সিজদা নয় (হানাফি) — যেমন রাফউল ইয়াদাইন।
\item \textbf{সন্দেহহীন, নিশ্চিত, সম্পূর্ণ} নামাজ হলে সাহু সিজদা লাগবে না।
\item একাকী \textbf{নফল} (যেমন ২-৪ রাকাত নফল) — ভুলের তীব্রতায় নিয়মে শিথিলতা আছে; তবে মূলনীতি বহাল।
\end{itemize}
\medskip\hrule\medskip

\subsection*{৬. ইমাম ও মাসবুকের জন্য সাহু সিজদা}

\begin{itemize}
\item \textbf{ইমাম ভুল করলে:} সাহু সিজদা \textbf{ইমাম ও মুক্তাদি সবাই} করে — ইমামই সাহু সিজদা করবেন, মুক্তাদি অনুসরণ করবেন (মুক্তাদির নিজের ভুল না থাকলে — তাকে সাহু সিজদা দিতে হবে না, ইমামের সাহু সিজদাই যথেষ্ট — হানাফি; জমহুরের সামান্য পার্থক্য)।
\item \textbf{মাসবুকের ক্ষেত্রে:}
\item মাসবুক \textbf{জামাতের রাকাতে} ইমামের ভুলে সাহু সিজদা হলে — সে জামাতি-রাকাতে অনুসরণ করবে; তার নিজের বাড়তি রাকাতে নিজের ভুলের জন্য নিজে সাহু সিজদা।
\item \textbf{সবচেয়ে কঠিন নিয়ম:} মাসবুকের সাহু সিজদা \textbf{তাঁর নিজের নামাজ শেষে} — ইমামের সালামের পর নিজের বাকি রাকাত পূর্ণ করে নিজের ভুল হলে। (বিস্তারিত উদাহরণসহ পার্ট ৭খ)
\end{itemize}
\medskip\hrule\medskip

\subsection*{৭. প্রশ্নোত্তর}

\begin{table}[h]\centering\small
\begin{tabular}{ll}
প্রশ্ন & উত্তর \\
\hline
ভুলে ৪-এর জায়গায় ৫ রাকাত পড়ে ফেললে? & সাহু সিজদা দিয়ে সংশোধন (বুখারি ৪০১) \\
সন্দেহ — ৩ না ৪? & নিশ্চিত ন্যূনতম (৩) ধরে নিন + সাহু সিজদা (মুসলিম ৫৭২) \\
সাহু সিজদা কখন? & সালামের আগে বা পরে — ভুলের ধরন ও মাযহাব অনুযায়ী \\
ইমামের ভুলে মুক্তাদি কি করবে? & ইমামের সাথে অনুসরণ করবে — নিজে অতিরিক্ত নয় (হানাফি) \\
সুন্নত ছুটলে সাহু সিজদা? & না — শুধু ফরজ/ওয়াজিব-সংক্রান্ত ভুলে \\
\hline
\end{tabular}
\end{table}

\medskip\hrule\medskip

\subsection*{৮. রেফারেন্স}

\begin{itemize}
\item সহীহ বুখারি — ৪০১, ৬৪২ (সন্দেহ-সংক্রান্ত ইমামের রুকন), ৭৫৭
\item সহীহ মুসলিম — ৫৭২, ৫৭৩
\item সুনানে আবু দাউদ — ১০১৯-১০২০ (সাহু সিজদার বিস্তারিত), ১০২৪
\item সুনানে ইবনে মাজাহ — ১২১১-১২১২ (সন্দেহের হাদিস)
\item আল-ফিকহুল ইসলামি ওয়া আদিল্লাতুহু (আল-যুহাইলি) — মাযহাব-অবস্থান সারণি
\end{itemize}
\begin{quote}
\textbf{দ্রষ্টব্য:} সাহু সিজদার সবচেয়ে জটিল কেস — যেমন রুকু ছুটে যাওয়া, সিজদা-রাকাতের সংখ্যা নিয়ে দ্বিধা, মাসবুকের সাহু সিজদা — পার্ট ৭খ (ভুল-ও-প্রতিকার) ফাইলে বিস্তারিত কেস-স্টাডি হিসেবে আছে।
\end{quote}

\end{document}
